\documentclass[a4paper,12pt]{article}

% Character set
\usepackage{cmap}
\usepackage[utf8]{inputenc}
\usepackage[T1]{fontenc} % ensure that all the characters in characterSets.tex prints
\usepackage{upquote} % \textcent
\usepackage{pifont} % add \ding, http://ctan.org/pkg/pifont

% A background text to prevent wide distribution
\usepackage{draftwatermark}
\SetWatermarkText{DRAFT}
\SetWatermarkScale{6}
\SetWatermarkLightness{.95}

% Page setup
\usepackage[top=25mm,bottom=20mm,inner=20mm,outer=40mm,marginparsep=3mm,marginparwidth=35mm]{geometry}
\renewcommand{\floatpagefraction}{.8}%

% paragraph indentation is stupid
\setlength\parindent{0pt}
\setlength{\parskip}{1em}

% Globally defined colors
\usepackage[table,x11names]{xcolor}
\definecolor{alternateKeywordsColor}{rgb}{0.13,1,0.13}
\definecolor{keywordsColor}{rgb}{0.13,0.13,1}
%\definecolor{commentsColor}{rgb}{0,0.5,0}
\definecolor{commentsColor}{rgb}{0,0.5,0}
%\definecolor{stringsColor}{rgb}{0.9,0,0}
\definecolor{stringsColor}{rgb}{0,0,0.5}
\definecolor{light-gray}{gray}{0.95}
\definecolor{codeLineHighlight}{named}{SlateGray1}
%\definecolor{codeLineHighlight}{rgb}{0.975,0.975,0.975}
\definecolor{syntaxColor}{rgb}{0,.45,0}

\definecolor{headerRowColor}{rgb}{0.85,0.85,0.85}
\definecolor{oddRowColor}{rgb}{0.95,0.95,0.95}
\definecolor{evenRowColor}{rgb}{1,1,1}

% add check- and crossmarks, http://ctan.org/pkg/pifont
\newcommand{\cmark}{{\color{green}\ding{51}}}%
\newcommand{\xmark}{{\color{red}\ding{55}}}%

% Extra math stuff
\usepackage{amsmath,amssymb}

% Typeset chess
\usepackage{skak}

% Figures
\usepackage{graphicx}
\graphicspath{{figures/}}

% clickable url
\usepackage{url}

% figures
\usepackage{subfigure}

% Clickable table of content
\usepackage[pdfpagelabels]{hyperref}
%\usepackage{multirow}
\usepackage{makecell}

% Include label name in ref
\usepackage[noabbrev,capitalize]{cleveref}
\newcommand{\creflastconjunction}{, and\nobreakspace~}
\Crefformat{tcb@cnt@codeNOutput}{Listing~#2#1#3}
\crefformat{tcb@cnt@codeNOutput}{Listing~#2#1#3}
\crefrangeformat{tcb@cnt@codeNOutput}{Listing~#3#1#4\nobreakdash--#5#2#6}
\Crefrangeformat{tcb@cnt@codeNOutput}{Listing~#3#1#4\nobreakdash--#5#2#6}
\crefmultiformat{tcb@cnt@codeNOutput}{Listing~#2#1#3}{ and~#2#1#3}{, #2#1#3}{\creflastconjunction#2#1#3}
\Crefmultiformat{tcb@cnt@codeNOutput}{Listing~#2#1#3}{ and~#2#1#3}{, #2#1#3}{\creflastconjunction#2#1#3}
\crefrangeformat{table}{Table~#3#1#4\nobreakdash--#5#2#6}
\Crefrangeformat{table}{Table~#3#1#4\nobreakdash--#5#2#6}
\crefrangeformat{part}{Part~#3#1#4\nobreakdash--#5#2#6}
\Crefrangeformat{part}{Part~#3#1#4\nobreakdash--#5#2#6}

% paragraphs in tables
\usepackage{tabularx}

% formatting lists
\usepackage{enumitem}
%\setlist[description]{leftmargin=0pt,labelindent=0pt,itemindent=0pt}
%\setlist[description]{itemindent=-\leftmargin}

% latex comment environment
\usepackage{comment}

% UML
\usepackage{pgf-umlcd}
\renewcommand{\umltextcolor}{black} 
\renewcommand{\umlfillcolor}{black!5!white}
\renewcommand{\umldrawcolor}{teal}

% List of indices
\usepackage{xstring}
\usepackage{makeidx}
\usepackage{marginfix} % fixes marginpar location problem in 2 -page mode.
\newcommand{\idxs}[1]{\marginpar{$\cdot$~\parbox[t]{\linewidth}{\raggedright \expandarg\IfSubStr{#1}{@}{\StrBehind{#1}{@}}{#1}}}\index{#1}} % The parbox is too wide, since the line also includes cdot-space
\newcommand{\idxss}[1]{\index{#1}}
% Define a new command idx with an optional parameter, which if given is the key to the index
\makeatletter
\def\idx{\@ifnextchar[{\@with}{\@without}}
\def\@with[#1]#2{\emph{#2}\idxs{#1}}
\def\@without#1{\emph{#1}\idxs{#1}}
\makeatother
%\newcommand{\idx}[1]{\emph{#1}\idxs{#1}}
\newcommand{\keyword}[1]{{\lstinline[language=fsharp]|#1|}}
\newcommand{\lexeme}[1]{\mbox{``\lstinline[language=fsharp]|#1|''}}
\makeindex

% display tree like structures
\usepackage{qtree}

% We frame all listings and problems
\usepackage{tcolorbox}
\tcbuselibrary{listings}
\tcbuselibrary{raster}
\tcbset{%
  colframe=teal, %PaleGreen1!45!black,
  %coltitle=black,
  fonttitle=\bfseries, 
  leftrule=3mm,
  sharp corners=downhill,
  colback=black!5!white,
  left=1mm,
  top=1mm,
  right=1mm,
  bottom=1mm,
  middle=1mm,
  arc=2mm,
}
\newtcolorbox[auto counter]{problem}[1][]{%
  title=\textbf{Problem~\thetcbcounter},
  colframe=DeepSkyBlue1, %green!30!blue,
  #1}
\newcommand{\src}{src}
\newtcolorbox[auto counter]{codeNOutput}[2][]{%
  title=\textbf{Listing~\thetcbcounter#2},
  #1}

%% lstlisting stuff
\usepackage{listings} 
\def\lstfs#1{\mbox{\lstinline{{#1}}}}
% Get counters from references for firstnumber references in lstinputlisting
\usepackage{refcount}
\newcounter{lstFrom}
\newcounter{lstTo}
% Example: 
% \setcounterref{lstFrom}{dynamicScopeTracing:a1}
% \setcounterref{lstTo}{dynamicScopeTracing:a2}
% \lstinputlisting[firstline=\thelstFrom,lastline=\thelstTo,escapechar=|]{\src/dynamicScopeTracing.fsx}
\usepackage{lstlinebgrd}
\makeatletter
%The following sets the box compatible with tcolorbox setup
\def\lst@linebgrdcolor{\color{black!5!white}}
\def\lst@linebgrdsep{1em}
\def\lst@linebackgroundwidth{1em}
\def\lst@linebackgroundhighlight{\color{codeLineHighlight}}
\renewcommand{\lst@linebgrd}{%
  \ifx\lst@linebgrdcolor\empty
  \else
    \rlap{
       \lst@basicstyle\color{black!5!white} % tcolorbox background color
       \lst@linebgrdcolor{
          \kern-\dimexpr\lst@linebgrdsep\relax
          \lst@linebgrdcmd{\lst@linebgrdwidth}{\lst@linebgrdheight}{\lst@linebgrddepth}
       }
    }
  \fi
}
% Highlight a range of lines with green. Use \getrefnumber{label} for refs
\newcommand{\highlightRange}[2]{\ifnum\value{lstnumber}>\numexpr#1-1\ifnum\value{lstnumber}<\numexpr1+#2\lst@linebackgroundhighlight\fi\fi}
% \highlight conflicts with skak. Just rewriting, wonder what breaks in skak
\renewcommand{\highlight}[1]{\ifnum\value{lstnumber}=#1\lst@linebackgroundhighlight\fi}

% To use verbatimwrite to write listing to file, e.g., in conjunction with ebnfs
\usepackage{moreverb} 

\lstdefinelanguage{fsharp}{%
  keywords={abstract, and, as, assert, base, begin, class, default, delegate, do, done, downcast, downto, elif, else, end, exception, extern, false, finally, for, fun, function, global, if, in, inherit, inline, interface, internal, lazy, let, match, member, module, mutable, namespace, new, null, of, open, or, override, private, public, rec, return, sig, static, struct, then, to, true, try, type, upcast, use, val, void, when, while, with, yield},
  morekeywords={atomic, break, checked, component, const, constraint, constructor, continue, eager, fixed, fori, functor, include, measure, method, mixin, object, parallel, params, process, protected, pure, recursive, sealed, tailcall, trait, virtual, volatile},
  otherkeywords={ let!, return!, do!, yield!, use!},
  keywordstyle=\color{keywordsColor},
  % sensitive=true,
  basicstyle=\ttfamily\lst@ifdisplaystyle\small\fi, % make font small for listings but not for lstinline
  breaklines=true,
  breakatwhitespace=true
  showstringspaces=false,
  morecomment=[l][\color{commentsColor}]{///},
  morecomment=[l][\color{commentsColor}]{//},
  morecomment=[n][\color{commentsColor}]{(*}{*)},
  morecomment=[is][\color{white}]{(*//}{*)},
  morestring=[b]",
  literate={`}{\`}1,
  stringstyle=\color{stringsColor},
  showspaces=true,
  numbers=left,
  numbersep=6pt,
  numberstyle=\scriptsize\color{white},
  % aboveskip=0pt, 
  % belowskip=0pt,
  % resetmargins=true,
  % captionpos=b,
  backgroundcolor=\color{black!5!white},
}


\lstdefinelanguage{syntax}{%
  classoffset=0,
  keywords={abstract, and, as, assert, base, begin, class, default, delegate, do, done, downcast, downto, elif, else, end, exception, extern, false, finally, for, fun, function, global, if, in, inherit, inline, interface, internal, lazy, let, match, member, module, mutable, namespace, new, null, of, open, or, override, private, public, rec, return, sig, static, struct, then, to, true, try, type, upcast, use, val, void, when, while, with, yield, atomic, break, checked, component, const, constraint, constructor, continue, eager, fixed, fori, functor, include, measure, method, mixin, object, parallel, params, process, protected, pure, recursive, sealed, tailcall, trait, virtual, volatile, let!, return!, do!, yield!, use!},
  keywordstyle=\color{keywordsColor},
  % classoffset=1,
  % morekeywords={ident, expr, arg, format-string},
  % keywordstyle=\color{syntaxColor},
  % classoffset=0,
  otherkeywords={},
  basicstyle=\ttfamily\lst@ifdisplaystyle\small\fi, % make font small for listings but not for lstinline
  breaklines=true,
  breakatwhitespace=true
  showstringspaces=false,
  classoffset=0,
  morecomment=[l][\color{commentsColor}]{////},
  literate={`}{\`}1 {\{*}{{{\color{syntaxColor}\{}}}1 {*\}}{{{\color{syntaxColor}\}}}}1 {[*}{{{\color{syntaxColor}[}}}1  {*]}{{{\color{syntaxColor}]}}}1 {|*}{{{\color{syntaxColor}|}}}1, % {etc*}{{{\color{syntaxColor}...}}}3,
  moredelim  = **[is][\processmydelims]{<*}{*>}, % delete delimiters, typeset keywords. Don't know how to avoid the last...
  showspaces=true,
  numbers=left,
  numbersep=6pt,
  numberstyle=\scriptsize\color{white},
  backgroundcolor=\color{black!5!white},
}
%Tweek of deliminter and literate: https://tex.stackexchange.com/questions/203263/listings-package-custom-language-delimiter-match-left-side
\newcommand\processmydelimsend{}
\newcommand\processmydelims{%
  \renewcommand\processmydelimsend{\textcolor{syntaxColor}{>}\egroup}%
  \bgroup\color{syntaxColor}<\aftergroup\processmydelimsend%
}
% \makeatletter
% \newcommand\processhash{%
%   \ifnum\lst@mode=\lst@Pmode%
%     \bfseries%
%   \fi
%   \#%
% }
% \makeatother


\lstdefinelanguage{ebnf}{%
  keywords={},
  morekeywords={},
  otherkeywords={},
  keywordstyle=\color{keywordsColor},
  % sensitive=true,
  basicstyle=\fontfamily{pcr}\selectfont\lst@ifdisplaystyle\small\fi, 
  breaklines=true,
  breakatwhitespace=true
  morecomment=[s][\color{commentsColor}]{(*}{*)},
  morestring=[b]",
  morestring=[b]',
  alsoletter={\\},
  showstringspaces=false,
  % stringstyle=\color{stringsColor},
  % aboveskip=0pt, 
  % belowskip=0pt,
  % resetmargins=true,
  % captionpos=b,
  % backgroundcolor=\color{blue!10!white},
}
\lstdefinelanguage{console}{%
  keywords={},
  morekeywords={},
  otherkeywords={},
  basicstyle=\ttfamily\lst@ifdisplaystyle\small\fi, 
  breaklines=true,
  showstringspaces=false,
  % aboveskip=0pt,
  % belowskip=0pt,
  % resetmargins=true,
  % captionpos=b,
  % backgroundcolor=\color{green!10!white},
}
%\lstset{language=fsharp, frame=single}
\lstset{language=fsharp,showlines=false}
\makeatletter
\def\lst@visiblespace{ }
\makeatother

% input .fsx and .out listings from \src and display as code and result in same figure
% #1 = optional further arguments for lstinputlisting
% #2 = filename without suffix, and label
% #3 = caption
\newtcbinputlisting[use counter from=codeNOutput]{\fs}[3][]{%
  listing file={src/#2.fsx},
  listing and comment,
  listing options={language=fsharp,escapechar=§,#1},
  title=\textbf{Listing \thetcbcounter} {#2.fsx:\\#3},
  label={#2},
  comment={\lstinputlisting[language=console]{\src/#2.out}}
}

% dispaly fsharp code \src
% #1 = optional further arguments for lstinputlisting
% #2 = filename
% #3 = label
% #4 = caption
\newtcbinputlisting[use counter from=codeNOutput]{\fsharp}[4][]{%
  listing file={\src/#2},
  listing only,
  listing options={language=fsharp,escapechar=§,#1},
  title=\textbf{Listing \thetcbcounter} {#2:\\#4},
  label={#3},
}

% dispaly console file \src
% #1 = optional further arguments for lstinputlisting
% #2 = filename
% #3 = label
% #4 = caption
\newtcbinputlisting[use counter from=codeNOutput]{\console}[4][]{%
  listing file={\src/#2},
  listing only,
  listing options={language=console,escapechar=§,#1},
  title=\textbf{Listing \thetcbcounter} {#2:\\#4},
  label={#3},
}

\newtcbinputlisting[use counter from=codeNOutput]{\fsCode}[4]{%
  listing file={src/#1.fsx},
  listing only,
  listing options={language=fsharp,escapechar=§,#4},
  title=\textbf{Listing \thetcbcounter} {#1.fsx:\\#3},
  label={#2},
}

% dispaly ebnf file, no label
% #1 = optional further arguments for lstinputlisting
% #2 = filename
% #3 = caption
\newtcbinputlisting[use counter from=codeNOutput]{\ebnf}[3][]{%
  listing file={#2},
  listing only,
  colframe=black!50!white,
  listing options={language=ebnf,escapechar=§,#1},
  title=\textbf{Listing \thetcbcounter} {#3},
}

% dispaly syntax file, no label
% #1 = optional further arguments for lstinputlisting
% #2 = filename without suffix, and label
% #3 = caption
\newtcbinputlisting[use counter from=codeNOutput]{\syntax}[3][]{%
  listing file={#2},
  listing only,
  colframe=black!50!white,
  listing options={language=syntax,escapechar=§,#1},
  title=\textbf{Listing \thetcbcounter} {#3},
  label={#2}
}

\newtcbinputlisting[use counter from=codeNOutput]{\fsSignature}[4]{%
  listing file={src/#1.fsi},
  listing only,
  listing options={language=fsharp,escapechar=§,#4},
  title=\textbf{Listing \thetcbcounter} {#1.fsi:\\#3},
  label={#2},
}
\newtcbinputlisting[use counter from=codeNOutput]{\fsImplementation}[4]{%
  listing file={src/#1.fs},
  listing only,
  listing options={language=fsharp,escapechar=§,#4},
  title=\textbf{Listing \thetcbcounter} {#1.fs:\\#3},
  label={#2},
}

% dispaly output file .out from \src
% #1 = optional further arguments for lstinputlisting
% #2 = filename without suffix, and label
% #3 = caption
\newtcbinputlisting[use counter from=codeNOutput]{\fsOutput}[3][]{%
  listing file={src/#2.out},
  listing only,
  listing options={language=console,escapechar=§,#1},
  title=\textbf{Listing \thetcbcounter}: {#3},
  label={#2},
}

% dispaly output file .out from \src as an element in tabularx
% #1 = optional further arguments for lstinputlisting
% #2 = filename without suffix, and label
% #3 = caption
\newtcbinputlisting[use counter from=codeNOutput]{\fsOutputTabx}[3][]{%
  listing file={src/#2.out},
  listing only,
  width=\hsize,
  box align=top,
  listing options={language=console,escapechar=§,aboveskip=0pt,belowskip=0pt,emptylines=0,#1},
  title=\textbf{Listing \thetcbcounter}: {#3},
  label={#2},
}

\newcommand{\filename}[1]{\lstinline[language=console]{#1}}

% highlighted text snippets
\newcommand{\advice}[1]{\marginpar{Advice}{\textbf{#1}}}
\newcommand{\advanced}[1]{\marginpar{Advanced material}\textbf{#1}}

% sometimes we need to include hash sign as arguments
\begingroup\catcode`\#=12
\newcommand\hashchar{}%check that is doesn't exist
\gdef\hashchar{#}
\endgroup

% Scratch out math, used in test.tex
\usepackage{cancel}
%\newcommand{\bcancel}[1]{#1}

% Draw arrows between elements
\usepackage{tikz}
%\usepackage{sphack} % make overlays invisible where stated in text
\usetikzlibrary{arrows,shapes,calc,decorations.pathreplacing}
\newcommand{\tikzmark}[1]{\tikz[overlay,remember picture] \node (#1) {};}
\newcommand*{\DrawArrow}[3][]{%
  % #1 = draw options
  % #2 = left point
  % #3 = right point
  \begin{tikzpicture}[overlay,remember picture]
    %\draw [-latex, #1,ultra thick,red] ($(#2)+(0.1em,0.5ex)$) to ($(#3)+(0,0.5ex)$);
    \draw [-latex, #1,ultra thick,red] ($(#2) -(0,0.5ex)$) to ($(#3)+(0,2ex)$);
  \end{tikzpicture}%
}%
\newcommand*{\AddNote}[4]{%
  \begin{tikzpicture}[overlay, remember picture]
    \draw [decoration={brace,amplitude=0.5em},decorate,ultra thick,red]
    ($(#3)!([yshift=1.5ex]#1)!($(#3)-(0,1)$)$) -- ($(#3)!(#2)!($(#3)-(0,1)$)$)
    node [align=left, text width=0cm, pos=0.5, anchor=west, xshift=.2cm] {#4};
  \end{tikzpicture}
}%
\newcommand{\FrameArea}[2]{%
  % #1 = top left point
  % #2 = bottom right point
  % The overlay is drawn in the margin in order not to screw with
  % horizontal spacing.
  %\ifvmode\vspace*{-1.2em}\else\fi%
  \begin{tikzpicture}[overlay,remember picture]%
    \draw[red,rounded corners] ([shift={(-2pt,1.9ex)}] #1)  rectangle  ([shift={(2pt,-.9ex)}] #2);%
  \end{tikzpicture}\noindent % I don't know why this command shift to the right, but this seems to fix it.
}%

% One can write to a file during compilation with the following
% low-level code.
%  \newwrite\tempfile
%  \immediate\openout\tempfile=list.tex
%  \immediate\write\tempfile{Text to write to file}
%  \immediate\closeout\tempfile

\usepackage{xspace}
\newcommand{\monoVersion}{5.2.0\xspace}
\newcommand{\fsharpVersion}{4.1\xspace}


% Notes to self
\newcommand{\jon}[1]{\footnote{Jon: \textbf{#1}}}
%\renewcommand{\jon}[1]{}
\newcommand{\mael}[1]{\footnote{Mael: \textbf{#1}}}
%\renewcommand{\mael}[1]{}
\newcommand{\spec}[1]{\footnote{Spec: \textbf{#1}}}
\renewcommand{\spec}[1]{}

%%% Local Variables:
%%% TeX-master: "fsharpNotes"
%%% End:

\newcommand{\id}[1]{\ensuremath{\emph{#1}}}
\title{Programmering og Problemløsning\\Datalogisk Institut,
  Københavns Universitet\\Arbejdsseddel 7 - individuel opgave}
\author{Martin Elsman}
\date{23.\ oktober -- 1.\ november.\\Afleveringsfrist: onsdag d.\ 1.\ november kl. 22:00}

\begin{document}
\maketitle

\noindent
I denne periode skal I arbejde individuelt. Regler for individuelle afleveringsopgaver er beskrevet i ``Noter, links, software m.m.'' $\rightarrow$ ''Generel information om opgaver''. Formålet er at arbejde med:
\begin{itemize}
%\item strenge
\item rekursion
\item pattern matching
\end{itemize}

Opgaverne er delt i øve- og afleveringsopgaver.

\section*{Øveopgaver}
\begin{enumerate}[label=5ø.\arabic*,start=0]
\item Omskriv funktionen \lstinline{insert}, som benyttes i forbindelse
  med funktionen \lstinline{isort} (insertion sort) fra forelæsningen,
  således at den benytter sig af pattern matching på lister.
\item Omskriv funktionen \texttt{bsort} (bubble sort) fra
  forelæsningen således at den benytter sig af pattern matching på
  lister. Funktionen kan passende benytte sig af ``nested pattern
  matching'' i den forstand at den kan implementeres med et match case
  der udtrækker de to første elementer af listen samt halen efter
  disse to elementer.
\item Opskriv black-box tests for de to sorteringsfunktioner og vær
  sikker på at grænsetilfældene dækkes (ingen elementer, et element,
  to elementer, samt flere elementer, sorteret, omvendt sorteret, etc.)

\item\label{sierpinski.ov} Ved at benytte biblioteket \texttt{ImgUtil}, som beskrevet i
  forelæsningen, er det muligt at tegne simpel liniegrafik samt
  fraktaler, som f.eks. Sierpinski-fraktalen, der kan tegnes ved at
  tegne små firkanter bestemt af et rekursivt mønster. Koden for
  Sierpinski-trekanten er givet som følger:

\begin{lstlisting}[numbers=none,frame=none,mathescape]
open ImgUtil

let rec triangle bmp len (x,y) =
  if len < 25 then setBox blue (x,y) (x+len,y+len) bmp
  else let half = len / 2
       do triangle bmp half (x+half/2,y)
       do triangle bmp half (x,y+half)
       do triangle bmp half (x+half,y+half)

do runSimpleApp "Sierpinski" 600 600 (fun bmp -> triangle bmp 512 (30,30) |> ignore)
\end{lstlisting}

Tilpas funktionen således at trekanten tegnes med røde streger samt
således at den kun tegnes ned til dybde 2 (hint: du skal ændre betingelsen \lstinline{len < 25}).

\item\label{sierpinskikeys.ov} I stedet for at benytte funktionen
  \lstinline{ImgUtil.runSimpleApp} skal du nu benytte funktionen
  \lstinline{ImgUtil.runApp}, som giver mulighed for at din løsning
  kan styres ved brug af tastaturet. Funktionen \lstinline{ImgUtil} har følgende type:
\begin{footnotesize}
\begin{lstlisting}[numbers=none,frame=none,mathescape]
val runApp : string -> int -> int
          -> (int -> int -> 's -> System.Drawing.Bitmap)
          -> ('s -> System.Windows.Forms.KeyEventArgs
                 -> 's option)
          -> 's -> unit
\end{lstlisting}
\end{footnotesize}
De tre første argumenter til \lstinline{runApp} er vinduets titel (en
streng) samt vinduets initielle vidde og højde. Funktionen
\lstinline{runApp} er parametrisk over en brugerdefineret type af tilstande (\lstinline{'s}). Antag at funktionen kaldes som følger:
\begin{footnotesize}
\begin{lstlisting}[numbers=none,frame=none,mathescape]
  runApp title width height draw react init
\end{lstlisting}
\end{footnotesize}
Dette kald vil starte en GUI applikation med titlen \lstinline{title},
vidden \lstinline{width} og højden \lstinline{height}. Funktionen
\lstinline{draw}, som brugeren giver som 4.\@ argument kaldes initielt
når applikationen starter og hver gang vinduets størrelse justeres
eller ved at funktionen \lstinline{react} er blevet kaldt efter en
tast er trykket på tastaturet. Funktionen \lstinline{draw} modtager
også (udover værdier for den aktuelle vidde og højde) en værdi for den
brugerdefinerede tilstand, som initielt er sat til værdien
\lstinline{init}. Funktionen skal returnere et bitmap, som for
eksempel kan konstrueres med funktionen \lstinline{ImgUtil.mk} og
ændres med andre funktioner i \lstinline{ImgUtil}
(f.eks. \lstinline{setPixel}).

Funktionen \lstinline{react}, som brugeren giver som 5.\@ argument
kaldes hver gang brugeren trykker på en tast. Funktionen tager som
argument en værdi svarende til den nuværende tilstand for
applikationen samt et argument der kan benyttes til at afgøre hvilken
tast der blev trykket på.\footnote{Hvis \lstinline{e} har typen
  \lstinline{System.Windows.Forms.KeyEventArgs} kan betingelsen
  \lstinline{e.KeyCode = System.Windows.Forms.Keys.Up} benyttes til at
  afgøre om det var tasten ``Up'' der blev trykket på.} Funktionen kan
nu (eventuelt) ændre på dens tilstand ved at returnere en ændret værdi
for denne.

Tilpas applikationen således at dybden af
fraktalen kan styres ved brug af piletasterne, repræsenteret ved værdierne
\lstinline{System.Windows.Forms.Keys.Up} og
\lstinline{System.Windows.Forms.Keys.Down}.
\end{enumerate}


\section*{Afleveringsopgaver}
\begin{enumerate}[label=5i.\arabic*,start=0]
\item Omskriv funktionen \lstinline{merge}, som benyttes i forbindelse
  med funktionen \lstinline{msort} (mergesort) fra forelæsningen,
  således at den benytter sig af pattern matching på lister.
\item Opskriv black-box tests for sorteringsfunktionen \lstinline{msort} og vær
  sikker på at grænsetilfældene dækkes (ingen elementer, et element,
  to elementer, samt flere elementer, sorteret, omvendt sorteret, etc.)
\end{enumerate}

\begin{minipage}{0.7\textwidth}
\begin{enumerate}[label=5i.\arabic*,start=2]
\item Med udgangspunkt i øvelsesopgave~\ref{sierpinski.ov} skal du i
  denne opgave implementere en GUI-applikation der kan tegne en
  version af X-fractalen som illustreret til højre (eventuelt i en
  dybde større end 2).

  Bemærk at det ikke er et krav at dybden på fraktalen skal kunne
  styres med piletasterne som det er tilfældet med
  Sierpinski-fraktalen i øvelsesopgave~\ref{sierpinskikeys.ov}.
\end{enumerate}
\end{minipage}\hspace{8mm} \begin{minipage}{0.25\textwidth}
  \includegraphics[width=0.9\textwidth]{xfractal.png}
\end{minipage}

\begin{enumerate}[label=5i.\arabic*,start=3]
\item Du skal i de følgende to opgaver arbejde med en funktion til at
  bestemme den såkaldte \emph{Levensthein-distance} mellem to strenge
  $a$ og $b$. Distancen er defineret som det mindste antal
  editeringer, på karakter-niveau, det er nødvendigt at foretage på
  strengen $a$ før den resulterende streng er identisk med strengen
  $b$. Som editeringer forstås (1) sletninger af karakterer, (2)
  indsættelser af karakterer, og (3) substitution af
  karakterer.

  Varianter af Levensthein-distancen mellem to strenge kan
  således benyttes til at identificere om studerende selv har løst
  deres indleverede opgaver eller om der potentielt set er tale om
  plagiatkode ;)

  Matematisk set kan Levensthein-distancen $\id{leven}(a,b)$, mellem
  to karakterstrenge $a$ og $b$, defineres som
  $\id{lev}_{a,b}(|a|,|b|)$, hvor $|a|$ og $|b|$ henviser til
  længderne af henholdsvis $a$ og $b$, og hvor funktionen $\id{lev}$ er
  defineret som følger:\footnote{See
    \url{https://en.wikipedia.org/wiki/Levenshtein_distance}.}

  \[
  \id{lev}_{a,b}(i,j) = \left \{ \begin{array}{ll} \mathrm{max}(i,j) & \mathrm{if}~\mathrm{min}(i,j) = 0, \\
    \mathrm{min} \left \{ \begin{array}{l}\id{lev}_{a,b}(i-1,j)+1 \\
                                          \id{lev}_{a,b}(i,j-1)+1 \\
                                          \id{lev}_{a,b}(i-1,j-1)+1_{(a_i\not=b_j)} \end{array} \right . & \mathrm{otherwise}. \end{array} \right .
  \]
  hvor $1_{(a_i\not=b_j)}$ henviser til \emph{indikatorfunktionen}, som er $1$ når $a_i\not=b_j$ og $0$ ellers.

  Implement\'{e}r funktionen \id{leven} direkte efter den matematiske
  definition (ved brug af rekursion) og test korrektheden af funktionen på nogle små
  strenge, såsom ``house'' og ``horse'' (distance 1) samt ``hi'' og ``hej'' (distance 2).

  \item Den direkte implementerede rekursive funktion er temmelig
    ineffektiv når strengene $a$ og $b$ er store. F.eks. tager det en
    del millisekunder at udregne distancen mellem strengene
    ``dangerous house'' and ``danger horse''. Årsagen til denne
    ineffektivitet er at en løsning der bygger direkte på den
    rekursive definition resulterer i en stor mængde genberegninger af
    resultater der allerede er beregnet.

    For at imødekomme dette problem skal du implementere en såkaldt
    ``caching mekanisme'' der har til formål at sørge for at en
    beregning højst foretages en gang.  Løsningen kan passende gøre
    brug af gensidig rekursion og tage udgangspunkt i løsningen for
    den direkte rekursive definition (således skal løsningen nu
    implementeres med to gensidigt rekursive funktioner \lstinline{leven} og \lstinline{leven_cache} forbundet med
    \lstinline{and}). Som cache skal der benyttes et 2-dimensionelt
    array af størrelse $|a|\times|b|$ indeholdende heltal (initielt
    sat til $-1$).

    Funktionen \lstinline{leven_cache}, der skal tage tilsvarende
    argumenter som \lstinline{leven}, skal nu undersøge om der
    allerede findes en beregnet værdi i cachen, i hvilket tilfælde
    denne værdi returneres. Ellers skal funktionen \lstinline{leven}
    kaldes og cachen opdateres med det beregnede resultat. Endelig er
    det nødvendigt at funktionen \lstinline{leven} opdateres til nu at
    kalde funktionen \lstinline{leven_cache} i hver af de rekursive kald.

    Test funktionen på de små strenge og vis at funktionen nu virker
    korrekt også for store input.

    Det skal til slut bemærkes at den implementerede løsning benytter
    sig af $O(|a| \times |b|)$ plads og at der findes effektive
    løsninger der benytter sig af mindre plads
    ($O(\mathrm{max}(|a|,|b|))$). Det er ikke et krav at din løsning
    implementerer en af disse mere pladsbesparende strategier.
\end{enumerate}

Afleveringsopgaven skal afleveres som et antal \texttt{fsx} tekstfiler
navngivet efter opgaven, som f.eks. \lstinline!7i0.fsx!. Tekstfilerne
skal kunne oversættes med \texttt{fsharpc},\footnote{Løsningerne kan
  antage at biblioteket \texttt{ImgUtil.dll} fra forelæsningerne er
  tilgængeligt.} og resultatet skal kunne køres med \texttt{mono}
eller eventuelt \texttt{mono32}. Funktioner skal dokumenteres ifølge
dokumentationsstandarden, og udover selve programteksten skal
besvarelserne indtastes som kommentarer i de \texttt{fsx}-filer, de
hører til. Det hele skal samles i en zip fil og uploades på Absalon.

Til øvelserne forventer vi at I arbejder efter følgende skema:
\begin{description}
\item[Mandag 23/10:] Afslut 6g og start på øvelsesopgaverne fra 7i
\item[Tirsdag 24/10:] Arbejd med øvelses- og afleveringsopgaverne
\item[Fredag 27/10] Arbejd med afleveringsopgaverne
\end{description}

\end{document}
